\section{Środowisko eksperymentalne} \label{sec:setup}

W tym rozdziale opisano sprzęt, konfigurację parametrów oraz talie użyte we wszystkich eksperymentach. Konfigurację tę przyjęto jako punkt odniesienia do porównań między wariantami.

\subsection{Sprzęt i oprogramowanie}

Wszystkie eksperymenty przeprowadzono na serwerze wyposażonym w dwa procesory AMD EPYC 7601 (łącznie 64 rdzenie fizyczne, 128 wątków logicznych), skonfigurowanymi jako jedna jednostka obliczeniowa, z 256 GB pamięci RAM DDR4 (2400 MHz) oraz systemem operacyjnym Windows Server 2019. Rozgrywki symulowano w środowisku SabberStone \cite{sabberstone_github}, który został napisany w języku C\# na platformie .NET Core. Skrypty optymalizacyjne uruchamiano w środowisku Python 3.

Rozgrywki nadają się do trywialnej równoległości (ang. \textit{embarrassingly parallel}) --- każda para agentów może być symulowana w osobnym wątku, co pozwala na wywoływanie walk jednocześnie na 128 wątkach procesora i pełne obciążenie dostępnych rdzeni. Pojedyncze uruchomienie eksperymentu trwało około 9,5 godziny. Wyjątek stanowi wariant z przeszukiwaniem głębokościowym (sekcja~\ref{sec:agent21depth}), który wymagał około 3,5 razy dłuższego czasu, tj. około 33 godzin, z uwagi na wzrost przestrzeni akcji przy głębszym przeszukiwaniu. Z tego powodu uruchamiany był tylko raz, aby wyrównać nakład obliczeniowy w stosunku do innych propozycji.

Łączny orientacyjny czas obliczeniowy wszystkich eksperymentów wyniósł około \textbf{572,81 godzin} ($\approx 23,87$ dni ciągłej pracy serwera). Zestawienie liczby uruchomień poszczególnych wariantów przedstawia tabela \ref{tab:runs}.

Dla większości wariantów wykonano po trzy niezależne przebiegi na każdy optymalizator. Liczba rzędu kilku powtórzeń jest typowa dla badań nad ewolucyjnym uczeniem funkcji oceny pozycji w grach, w których wysoki koszt obliczeniowy pojedynczego przebiegu wymusza ograniczenie liczby powtórzeń — przykładowo, Jaśkowski i Szubert w analogicznym zadaniu uczenia wag funkcji oceny dla gry Othello powtarzali każdą konfigurację pięciokrotnie \cite{jaskowski2016}. Przyjęte w niniejszej pracy trzy powtórzenia dla większości konfiguracji mieszczą się zatem w praktyce tego typu eksperymentów.

\begin{table}[p]
\centering
\small
\caption{Zestawienie uruchomień eksperymentów.}
\label{tab:runs}
\begin{tabular}{llcc}
\toprule
Wariant & Schemat oceny & Uruch. & Czas [h] \\
\midrule
\multicolumn{4}{l}{\textit{Modyfikacje agenta bazowego (koewolucja, POP\_SIZE = 10)}} \\
\hspace{1em}\textit{28 parametrów}      & koewolucja & 3 & 27,71 \\
\hspace{1em}\textit{28 znormalizowany}  & koewolucja & 3 & 29,71 \\
\hspace{1em}\textit{63 parametry}       & koewolucja & 3 & 27,09 \\
\hspace{1em}\textit{63 płynny}          & koewolucja & 3 & 26,86 \\
\hspace{1em}\textit{21 głębokościowy}   & koewolucja & 1 & 34{,}51 \\
\midrule
\multicolumn{4}{l}{\textit{Tuning parametru POP\_SIZE dla SHADE}} \\
\hspace{1em}SHADE koew., POP\_SIZE = 5    & koewolucja   & 2 & 18,81 \\
\hspace{1em}SHADE koew., POP\_SIZE = 10   & koewolucja   & 2 & 18,14 \\
\hspace{1em}SHADE koew., POP\_SIZE = 15   & koewolucja   & 2 & 17,83 \\
\hspace{1em}SHADE HoF, POP\_SIZE = 5      & Hall of Fame & 2 & 18,86 \\
\hspace{1em}SHADE HoF, POP\_SIZE = 10     & Hall of Fame & 2 & 19,31 \\
\hspace{1em}SHADE HoF, POP\_SIZE = 15     & Hall of Fame & 2 & 19,37 \\
\midrule
\multicolumn{4}{l}{\textit{Testy SHADE na agencie bazowy 21}} \\
\hspace{1em}\textit{bazowy 21}        & SHADE koew.  & 1 & 8{,}89 \\
\hspace{1em}\textit{bazowy 21}        & SHADE HoF    & 1 & 9{,}70 \\
\midrule
\multicolumn{4}{l}{\textit{Warianty SHADE koewolucyjnego (POP\_SIZE = 15)}} \\
\hspace{1em}\textit{28 parametrów}      & SHADE koew. & 3 & 27,81 \\
\hspace{1em}\textit{28 znormalizowany}  & SHADE koew. & 3 & 29,37 \\
\hspace{1em}\textit{63 parametry}       & SHADE koew. & 3 & 26,57 \\
\hspace{1em}\textit{63 płynny}          & SHADE koew. & 3 & 26,42 \\
\hspace{1em}\textit{21 głębokościowy}   & SHADE koew. & 1 & 30{,}21 \\
\midrule
\multicolumn{4}{l}{\textit{Warianty SHADE z Hall of Fame (POP\_SIZE = 15)}} \\
\hspace{1em}\textit{28 parametrów}      & SHADE HoF & 3 & 30{,}36 \\
\hspace{1em}\textit{28 znormalizowany}  & SHADE HoF & 3 & 31{,}78 \\
\hspace{1em}\textit{63 parametry}       & SHADE HoF & 3 & 29{,}41 \\
\hspace{1em}\textit{63 płynny}          & SHADE HoF & 3 & 29{,}47 \\
\hspace{1em}\textit{21 głębokościowy}   & SHADE HoF & 1 & 34{,}62 \\
\midrule
\textbf{Suma} & & \textbf{53} & \textbf{572,81} \\
\bottomrule
\end{tabular}
\end{table}

\subsection{Dobór rozmiaru populacji} \label{sec:popsize_tuning}

Parametry strategii ewolucyjnej w wariantach koewolucyjnych przyjęto bezpośrednio z badania García-Sáncheza i in.~\cite{sanchez2020} --- autorzy przeprowadzili własny dobór konfiguracji dla agenta bazowego. Dla metody SHADE sytuacja jest inna, ponieważ wśród jego parametrów ważny wpływ na jakość treningu i koszt obliczeniowy ma rozmiar populacji \texttt{POP\_SIZE}. Z tego powodu przed właściwymi eksperymentami przeprowadzono strojenie tego parametru. Aby nie faworyzować SHADE dodatkowym strojeniem na każdym wariancie agenta, testowano wyłącznie agenta bazowego (21 parametrów) --- uzyskana wartość \texttt{POP\_SIZE} jest następnie stosowana we wszystkich wariantach SHADE w trakcie badań.

Przetestowano wartości $\texttt{POP\_SIZE} \in \{5, 10, 15\}$ dla obu wariantów SHADE (koewolucyjnego i z Hall of Fame). Dla każdej wartości wytrenowano dwóch niezależnych agentów, uzyskując dwa różne zestawy wag (w tabeli~\ref{tab:popsize}: zestaw~1 i zestaw~2). Łącznie powstało więc $3 \cdot 2 = 6$ konfiguracji na wariant SHADE. W turnieju porównawczym zestawiono wszystkich sześciu agentów w formacie każdy z każdym (270 000 gier, po 2000 na parę kombinacji talii), osobno dla każdego wariantu SHADE.

Wyniki turnieju zestawiono w tabeli~\ref{tab:popsize}. W obu wariantach agenci wytrenowani z $\texttt{POP\_SIZE} = 15$ uzyskali najwyższy wskaźnik wygranych --- odpowiednio 52,15\% i 51,65\% (wariant koewolucyjny) oraz 52,30\% i 51,20\% (wariant Hall of Fame), wyraźnie przewyższając konfiguracje z $\texttt{POP\_SIZE} = 5$, które osiągnęły ok.\ 48\%--50\%. Na podstawie tych wyników przyjęto $\texttt{POP\_SIZE} = 15$ jako wartość bazową dla dalszych eksperymentów z użyciem SHADE.

\begin{table}[!h]
\centering
\caption{Wyniki turnieju doboru rozmiaru populacji}  \label{tab:popsize}
\begin{tabular}{llrrr}
\toprule
Wariant & Agent & Wygrane & Gry & Zwycięstwa\% \\
\midrule
\multirow{6}{*}{Koewolucyjny}
    & POP\_SIZE = 5,  zestaw 1 & 43\,775 & 90\,000 & 48,64\% \\
    & POP\_SIZE = 5,  zestaw 2 & 43\,312 & 90\,000 & 48,12\% \\
    & POP\_SIZE = 10, zestaw 1 & 45\,824 & 90\,000 & 50,92\% \\
    & POP\_SIZE = 10, zestaw 2 & 43\,663 & 90\,000 & 48,51\% \\
    & \textbf{POP\_SIZE = 15, zestaw 1} & \textbf{46\,938} & \textbf{90\,000} & \textbf{52,15\%} \\
    & \textbf{POP\_SIZE = 15, zestaw 2} & \textbf{46\,488} & \textbf{90\,000} & \textbf{51,65\%} \\
\midrule
\multirow{6}{*}{Hall of Fame}
    & POP\_SIZE = 5,  zestaw 1 & 45\,150 & 90\,000 & 50,17\% \\
    & POP\_SIZE = 5,  zestaw 2 & 40\,291 & 90\,000 & 44,77\% \\
    & POP\_SIZE = 10, zestaw 1 & 46\,232 & 90\,000 & 51,37\% \\
    & POP\_SIZE = 10, zestaw 2 & 45\,180 & 90\,000 & 50,20\% \\
    & \textbf{POP\_SIZE = 15, zestaw 1} & \textbf{46\,079} & \textbf{90\,000} & \textbf{51,20\%} \\
    & \textbf{POP\_SIZE = 15, zestaw 2} & \textbf{47\,068} & \textbf{90\,000} & \textbf{52,30\%} \\
\bottomrule
\end{tabular}
\end{table}

\subsection{Parametry algorytmów}
\label{sec:params}

Parametry strategii ewolucyjnej dla wariantów koewolucyjnych przyjęto zgodnie z konfiguracją oryginalnego badania~\cite{sanchez2020}, aby zachować porównywalność wyników i zbadać wyłącznie wpływ samych modyfikacji funkcji oceny. Jedyną zmianą względem oryginału jest zwiększenie parametru odpowiedzialnego za liczbę rdzeni, co wynika  z dostępności wieloprocesorowego serwera. Parametry te zestawiono w tabeli~\ref{tab:params_base}.

\begin{table}[ht]
\centering
\caption{Parametry dla modyfikacji agenta bazowego (warianty koewolucyjne).}
\label{tab:params_base}
\begin{tabular}{lll}
\toprule
Parametr & Wartość & Opis \\
\midrule
\texttt{NUM\_GAMES}       & 20           & Liczba partii na parę / na kombinację talii \\
\texttt{POP\_SIZE}        & 10           & Rozmiar populacji \\
\texttt{MAX\_EVALUATIONS} & 1000         & Maksymalna liczba ewaluacji osobników \\
\texttt{NUM\_THREADS}     & 128          & Liczba równoległych wątków symulacji \\
$D$                       & 21 / 28 / 63 & Wymiar przestrzeni wag \\
\bottomrule
\end{tabular}
\end{table}

Warianty SHADE używają większej populacji ($\texttt{POP\_SIZE} = 15$) dobranej w sekcji~\ref{sec:popsize_tuning} na agencie bazowym. Maksymalne okno historii parametrów SHADE ustawiono na $H = 100$, jednak
przy $\texttt{POP\_SIZE} = 15$ historia zawiera co najwyżej:
\[
\left\lfloor \frac{\texttt{MAX\_EVALUATIONS}}{\texttt{POP\_SIZE}} \right\rfloor
= \left\lfloor \frac{645}{15} \right\rfloor = 43
\]
dla wariantu koewolucyjnego oraz $\lfloor 525 / 15 \rfloor = 35$ dla wariantu
z Hall of Fame, więc pełna pojemność bufora nie jest wykorzystywana.. Wartości \texttt{MAX\_EVALUATIONS} różnią się między wariantami, co wynika z faktu, że SHADE z Hall of Fame ma mniejszy budżet ewaluacji. W tym wariancie co $K_{\text{HOF}} = 5$ generacji wszyscy rodzice są ponownie ewaluowani po aktualizacji $\mathcal{R}$, co generuje dodatkowy koszt obliczeniowy nieujęty w liczniku \texttt{MAX\_EVALUATIONS}. Parametry obu wariantów SHADE zestawiono odpowiednio w tabelach \ref{tab:params_coev} oraz \ref{tab:params_hof}.

\begin{table}[ht]
\centering
\caption{Parametry dla SHADE koewolucyjnego.}
\label{tab:params_coev}
\begin{tabular}{lll}
\toprule
Parametr & Wartość & Opis \\
\midrule
\texttt{NUM\_GAMES}       & 20           & Liczba partii na parę / na kombinację talii \\
\texttt{POP\_SIZE}        & 15           & Rozmiar populacji (po procesie dostrajania) \\
\texttt{MAX\_EVALUATIONS} & 645          & Maksymalna liczba ewaluacji osobników \\
\texttt{NUM\_THREADS}     & 128          & Liczba równoległych wątków symulacji \\
$H$                       & 100          & Maksymalne okno historii parametrów SHADE \\
$D$                       & 21 / 28 / 63 & Wymiar przestrzeni wag \\
\bottomrule
\end{tabular}
\end{table}

\begin{table}[ht]
\centering
\caption{Parametry dla SHADE z Hall of Fame.}
\label{tab:params_hof}
\begin{tabular}{lll}
\toprule
Parametr & Wartość & Opis \\
\midrule
\texttt{NUM\_GAMES}          & 20           & Liczba partii na parę / na kombinację talii \\
\texttt{POP\_SIZE}           & 15           & Rozmiar populacji (po procesie dostrajania) \\
\texttt{MAX\_EVALUATIONS}    & 525          & Maksymalna liczba ewaluacji osobników \\
\texttt{NUM\_THREADS}        & 128          & Liczba równoległych wątków symulacji \\
$H$                          & 100          & Maksymalne okno historii parametrów SHADE \\
$|\mathcal{R}|$              & 29           & Liczba agentów w pliku referencyjnym \\
$K_{\text{HOF}}$             & 5            & Interwał aktualizacji Hall of Fame [gen.] \\
$D$                          & 21 / 28 / 63 & Wymiar przestrzeni wag \\
\bottomrule
\end{tabular}
\end{table}


\subsection{Talie eksperymentalne}

Eksperymenty przeprowadzono z użyciem trzech talii stosowanych w oryginalnym badaniu \cite{sanchez2020}, a każda para agentów rozgrywa partie we wszystkich dziewięciu kombinacjach talii ($3 \cdot 3$), co zapewnia ocenę odporną na specyfikę konkretnej talii. Pełny skład talii przedstawia tabela \ref{tab:decklists}.

\paragraph{\textit{RenoKazakusMage} --- klasa Mag.}
Jest to talia skupiająca uwagę na kontroli planszy i opiera się na unikalnej mechanice Kazakusa --- każda karta w talii jest unikalna, nie ma swoich duplikatów, co aktywuje zdolność Reno Jacksona, czyli pełne leczenie bohatera. Strategia polega na przeżyciu wczesnej gry, uniemożliwieniu przeciwnikowi rozbudowy planszy i wygraniu poprzez silne zaklęcia w późnej fazie. Talia faworyzuje agentów zarządzających zasobami kosztem agresywnych strategii.

\paragraph{\textit{MidrangeJadeShaman} --- klasa Szaman.}
Jest to talia, która stara się balansować rozgrywkę i opiera się na mechanice Jade Golem --- $n$-ty przywołany golem ma $n$ punktów ataku oraz $n$ punktów zdrowia, więc każdy kolejny jest o jeden punkt ataku i o jeden punkt zdrowia silniejszy niż poprzedni. Strategia łączy wczesną kontrolę planszy z narastającą przewagą w fazie środkowej i końcowej. Mechanika ta wymaga od agenta uwzględnienia wartości stronników w horyzoncie wielu tur, a nie tylko aktualnego stanu planszy.

\paragraph{\textit{AggroPirateWarrior} --- klasa Wojownik.}
Talia agresywna oparta na mechanice piratów i synergii z bronią bohatera. Celem jest zadanie jak największych obrażeń bezpośrednio bohaterowi przeciwnika, zanim agresja opadnie i przewaga zasobów przejdzie na stronę przeciwnika. Talia premiuje agentów rozumiejących wartość broni i synergii z piratami.

\begin{table}[p]
\centering
\small
\caption{Składy talii użytych w eksperymentach. Karty posegregowane alfabetycznie. Liczba przy nazwie oznacza liczbę kopii w talii.}
\label{tab:decklists}
\begin{tabular}{p{4.8cm} p{4.8cm} p{4.8cm}}
\toprule
\textbf{\textit{AggroPirateWarrior}} & \textbf{\textit{MidrangeJadeShaman}} & \textbf{\textit{RenoKazakusMage}} \\
\textbf{(klasa Wojownik)} & \textbf{(klasa Szaman)} & \textbf{(klasa Mag)} \\
\midrule
Arkanitowy Rozpruwacz ×2        & Al'Akir, Władca Wichrów ×1   & Alexstraza ×1 \\
Elita Kor'kronu ×2              & Aya Czarny Pazur ×1          & Bełkocząca Księga ×1 \\
Grabieżca Krwawych Żagli ×2     & Barnes ×1                    & Brann Miedziobrody ×1 \\
Heroiczne Cios ×2               & Błyskawica ×1                & Brudny Szczur ×1 \\
Kapitan Południowych Mórz ×2    & Burza z Piorunami ×1         & Fala Płomieni ×1 \\
Kaprawe Oczko ×1                & Istota z Otchłani ×2         & Kozakus ×1 \\
Korsarz Naga ×2                 & Kaprawe Oczko ×1             & Księga Kabalisty ×1 \\
Kultystka Krwawych Żagli ×2     & Lazurowy Smokowiec ×2        & Kula Ognia ×1 \\
Majtek z Południowych Mórz ×2   & Mag Krwi Thalnos ×1          & Kurierka Konfraterni ×1 \\
Ognisty Wojenny Topór ×2        & Nefrytowa Błyskawica ×2      & Lazurowy Smokowiec ×1 \\
Pierwszy Mat N'Zotha ×2         & Nefrytowe Pazury ×2          & Lodowa Bariera ×1 \\
Pomniejszy Bukanier ×2          & Pomniejszy Bukanier ×2       & Lodowy Blok ×1 \\
Rozjuszony Berserker ×2         & Portal Maelstromu ×2         & Mag Krwi Thalnos ×1 \\
Sir Płetwin Mrrgglton ×1        & Ragnaros, Władca Ognia ×1    & Pocisk Lodu ×1 \\
Straszliwy Korsarz ×2           & Totemiczny Golem ×2          & Polimorfia ×1 \\
Ulepszenie! ×2                  & Totem Ognistego Języka ×2    & Portal Krainy Ognia ×1 \\
                                & Trogg Tunelarz ×2            & Prorok Zagłady ×1 \\
                                & Widmowe Pazury ×2            & Reno Jackson ×1 \\
                                & Zły Urok ×2                  & Sprzedawca Przekąsek ×1 \\
                                &                              & Sylwana Bieżywiatr ×1 \\
                                &                              & Szalona Duszomantka ×1 \\
                                &                              & Tajemny Intelekt ×1 \\
                                &                              & Tajemny Wybuch ×1 \\
                                &                              & Technik Kontroli Umysłu ×1 \\
                                &                              & Wulkaniczna Mikstura ×1 \\
                                &                              & Zakazany Płomień ×1 \\
                                &                              & Zamieć ×1 \\
                                &                              & Zapomniana Pochodnia ×1 \\
                                &                              & Żrący Bagienny Szlam ×1 \\
                                &                              & Żywiołak Wody ×1 \\
\bottomrule
\end{tabular}
\end{table}
